\chapter{Conclusiones}\label{conclusiones}

El principal objetivo de este proyecto es acercar el protocolo diseñado por de Feo, Jao y Plût a una persona con un leve conocimiento de criptografía. Para ello, hemos realizado una investigación sobre las bases matemáticas del protocolo que después hemos simplificado en el capítulo \ref{conceptosprevios} titulado Conceptos previos. Esta investigación ha sido extensa, ya que para entender íntegramente la extensa base que tiene SIKE hemos dado, a lo mejor, unos cuantos rodeos de más. No obstante, estos han sido finalmente necesarios para determinar qué conceptos eran realmente importantes por sus implicaciones y cuales no.

Naturalmente, este trabajo incluye también la redacción y revisión continua de los capítulos didácticos que son el capítulo \ref{conceptosprevios}  ya mencionado así como el capítulo \ref{estudioteorico} titulado Estudio teórico. 

De manera paralela se ha realizado una implementación gráfica que a nivel de interfaz refleja lo redactado en el capítulo \ref{estudioteorico} y a nivel de código ejemplifica cómo utilizar la implementación de SIKE en Java.

En nuestras investigaciones, tanto sobre el protocolo en sí como búsquedas técnicas específicas de problemas de desarrollo, sólo hemos encontrado un otro trabajo parecido a este. Se trata de una tesis equivalente a un trabajo fin de grado titulado \textit{Implementation of SIKE in Java: A detailed and comprehensive guide for understanding and implementing SIKE - A NIST Post-Quantum-Competition candidate} \cite{sike-implementation}. En cambio, este trabajo se centra en buscar una optimización de SIKE. Se desarrolla una implementación original de SIDH en SageMath y una de SIDH/SIKE en Java y, aunque explica las funciones más importantes, lo hace de manera formal. Este trabajo por Schubert no está pensado para iniciar a alguien a SIKE sino a desarrollar SIKE como protocolo.

En relación al desarrollo del proyecto, me siento orgullosa de la labor de investigación realizada. Se tratan temas difíciles y aunque ha habido veces que he sentido que estaba entrando en una agujero de conejo, una vez terminado he disfrutado la experiencia de desmigar un problema, encontrar las bases y construir mi propio conocimiento a partir de ahí.

De igual manera, el desarrollo de la interfaz ha sido bastante manual. Acostumbrada a empezar de una plantilla para todo proyecto que hemos realizado en clase (por motivos de tiempo y comodidad, principalmente) me propuse hacer este proyecto, entender realmente como se conecta todo. Hoy en día que la inteligencia artificial parece ser capaz de hacerlo todo en segundos, he querido centrarme en las bases que nunca antes había aprendido, 

Como he mencionado, el desarrollo es casi íntegramente personal. Usando la inteligencia artificial (el agente por defecto de Cursor) para corregir los fallos del entorno de pruebas del frontend, ya que llevaba días intentando configurarlo pero no conseguía que se llegasen siquiera a ejecutar. El resto des trabajo, investigación, redacción y desarrollo, es propio.

Este trabajo ha avanzado más lento de lo que tenía previsto, debido a la complejidad de las matemáticas pero también por cuestiones personales. Aun así, estoy contenta con el trabajo realizado. Habiendo empezado el doble grado en matemáticas e informática, decidí hacer este TFG, probablemente más matemático de lo normal en esta carrera, por volver a acercarme a esta disciplina. 